Supervised contrastive learning can use labels during pretraining, but in practice labels are scarce and
unlabeled images are plentiful. Should the few labels go inside the contrastive loss, or be kept until
after label-free pretraining? Two confounds hide the answer: unequal compute, and evaluations that use labels
the practitioner does not have. We remove both with a compute-matched $2\times2$ design: SimCLR and supervised
contrastive learning (SupCon), each trained on the labeled subset or on all images for the same number of
gradient steps, on CIFAR-10, CIFAR-100, STL-10 and Tiny-ImageNet. Three results stand out. First, labels inside
the loss buy speed, not a better end point: with 10\% of the CIFAR-10 labels, SupCon leads SimCLR by
\R{d:simclr-supcon:lp50:e10}\pt after 970 steps and trails it by \R{long:d:simclr-supcon}\pt after
\R{long:steps}. Second, the protocol and the use of the labels change the comparison: after 9,700 steps the
standard all-label linear probe puts a semi-supervised SupCon \R{cmp:simclr-union:lp50:e100:diff}\pt
ahead of SimCLR, fine-tuning with the default recipe puts it \R{cmp:simclr-union:ft:e100:diff}\pt behind when all 5,000 labels went into
pretraining, and level when 1,000 are held out to choose the recipe; with longer training SimCLR leads under every
probe and recipe. Third, a seen/unseen
diagnostic shows that label-using encoders make classes several times tighter on their labeled images than on any
other images. Given enough pretraining, scarce labels are better kept out of the contrastive loss and spent
afterwards, best of all in semi-supervised training. Budgets should be counted in steps, and encoders evaluated with the available
labels.
